\documentclass[12pt,a4paper]{article}
\usepackage[utf8]{inputenc}
\usepackage[T1]{fontenc}
\usepackage[french]{babel}
\usepackage{csquotes}
\usepackage[margin=2.5cm]{geometry}
\usepackage{amsmath,amssymb,amsthm}
\usepackage{setspace}
\usepackage{microtype}
\usepackage{xcolor}
\usepackage{hyperref}
\hypersetup{colorlinks=true,linkcolor=blue!50!black,urlcolor=blue!50!black,citecolor=blue!50!black}
\onehalfspacing

\newtheorem{theoreme}{Théorème}
\newtheorem{lemme}[theoreme]{Lemme}
\newtheorem{proposition}[theoreme]{Proposition}
\newtheorem*{remarque}{Remarque}

\title{\textbf{De l'unique identité de Catalan dans la famille des radicaux emboîtés de Ramanujan}}
\author{Yan Senez\\\texttt{yan.senez@protonmail.com}}
\date{\today}

\begin{document}
\maketitle

\begin{abstract}
Pour tout entier $n \geq 1$, l'identité de Ramanujan en radical emboîté
\[
n + 1 \;=\; \sqrt{\,1 + n\sqrt{1 + (n+1)\sqrt{1 + (n+2)\sqrt{\cdots}}}\,}
\]
est vérifiée, avec déploiement de premier niveau $(n+1)^2 = 1 + n(n+2)$. Nous
montrons que le cas $n = 2$, donnant $3^2 - 2^3 = 1$, est l'unique membre de
la famille pour lequel la quantité $n(n+2)$ est une puissance parfaite non
triviale. La preuve se réduit à deux observations élémentaires et à une
application du théorème de Mihailescu~\cite{Mihailescu2004}, qui résout la
conjecture de Catalan. Nous commentons ce que cela dit de la place du cas
$n = 2$ au sein de la famille de Ramanujan.
\end{abstract}

\section*{1. Introduction}

Dans un problème soumis au \emph{Journal of the Indian Mathematical
Society}~\cite{Ramanujan1911}, Ramanujan demande au lecteur d'évaluer
\begin{equation}
\sqrt{\,1 + 2\sqrt{1 + 3\sqrt{1 + 4\sqrt{1 + 5\sqrt{1 + \cdots}}}}\,},
\label{eq:Rama-3}
\end{equation}
donnant la réponse $3$. Ce radical emboîté est l'une des identités les plus
célèbres de l'arithmétique élémentaire. C'est en fait un cas particulier de
la formule générale consignée dans les \emph{Notebooks} de
Ramanujan~\cite[chap.~12]{Berndt1985}~:
\begin{equation}
n + 1 \;=\; \sqrt{\,1 + n\sqrt{1 + (n+1)\sqrt{1 + (n+2)\sqrt{\cdots}}}\,},
\qquad n \in \mathbb{N}_{\geq 1}.
\label{eq:Rama-general}
\end{equation}
La convergence du membre de droite, vue comme limite de troncatures finies,
a été établie par Vijayaraghavan~\cite{Vijayaraghavan1929} et
Herschfeld~\cite{Herschfeld1935}~; l'identité algébrique elle-même découle de
la relation récursive $f(n)^2 = 1 + n \cdot f(n+1)$ avec $f(n) = n+1$, puisque
$(n+1)^2 - 1 = n(n+2)$.

Le cas $n = 2$ dans~(\ref{eq:Rama-general}) redonne~(\ref{eq:Rama-3}). Au-delà
de son élégance visuelle, il est naturel de se demander si $n = 2$ jouit
d'une singularité \emph{intrinsèque} au sein de la famille
(\ref{eq:Rama-general}). À première vue, non~: l'identité est paramétrique,
et tout $n \geq 1$ produit un radical emboîté parfaitement valide. Dans cette
note, nous observons que, lorsque l'on regarde le développement au premier
niveau
\begin{equation}
(n+1)^2 \;=\; 1 + n(n+2),
\label{eq:first-level}
\end{equation}
le cas $n = 2$ jouit bel et bien d'une propriété singulière~: c'est l'\emph{unique}
valeur de $n$ pour laquelle $n(n+2)$ est une puissance parfaite non triviale,
et par conséquent l'unique valeur pour laquelle (\ref{eq:first-level}) est
elle-même une identité de type Catalan
\begin{equation}
3^2 - 2^3 \;=\; 1.
\label{eq:Catalan-3-2}
\end{equation}

Le cœur de l'observation est que cette unicité n'est pas seulement empirique~:
elle est une conséquence du théorème de Mihailescu~\cite{Mihailescu2004}, la
preuve de la conjecture de Catalan. L'objet de cette note est de consigner
cette remarque, d'en donner une preuve propre, et de discuter ce que cela
dit de la place de $n = 2$ dans la famille de Ramanujan.

\section*{2. Le résultat}

\begin{theoreme}\label{thm:main}
Pour $n \in \mathbb{Z}_{\geq 1}$, l'entier $n(n+2)$ est une puissance
parfaite non triviale, c'est-à-dire qu'il existe $b, q \in \mathbb{Z}_{\geq 2}$
tels que
\[
n(n+2) \;=\; b^{q},
\]
si et seulement si $n = 2$, auquel cas $(b, q) = (2, 3)$ et
l'identité~(\ref{eq:first-level}) se réduit à $3^{2} - 2^{3} = 1$.
\end{theoreme}

\begin{proof}
Nous distinguons selon la valeur de $q \geq 2$.

\medskip
\noindent\textbf{Cas 1~: $q = 2$.} Supposons $n(n+2) = b^2$ pour un certain
$b \geq 2$. Avec $n(n+2) = (n+1)^2 - 1$, ceci équivaut à
\begin{equation}
(n+1)^2 - b^2 \;=\; 1,
\quad\text{i.e.,}\quad
\bigl((n+1) - b\bigr)\bigl((n+1) + b\bigr) \;=\; 1.
\label{eq:diff-squares}
\end{equation}
Comme $b \geq 2$ et $n \geq 1$, les deux facteurs du membre de gauche
de~(\ref{eq:diff-squares}) sont des entiers strictement positifs~; or
l'unique factorisation de $1$ dans $\mathbb{Z}_{>0}$ est $1 \cdot 1$, ce qui
impose $(n+1) - b = (n+1) + b = 1$, donc $b = 0$ et $n + 1 = 1$, en
contradiction avec $b \geq 2$ et $n \geq 1$. Ainsi aucun $n \geq 1$ ne rend
$n(n+2)$ carré parfait.

(De manière équivalente~: $(n+1)^2 - 1$ se situe strictement entre les
carrés consécutifs $n^2$ et $(n+1)^2$ pour $n \geq 1$, donc n'est pas un
carré.)

\medskip
\noindent\textbf{Cas 2~: $q \geq 3$.} Supposons $n(n+2) = b^q$ pour un
$b \geq 2$ et un $q \geq 3$. Alors
\begin{equation}
(n+1)^2 - b^q \;=\; 1,
\label{eq:cat-shape}
\end{equation}
ce qui est l'équation de Catalan $a^p - b^q = 1$ avec les paramètres
$(a, p, q) = (n+1, 2, q)$. D'après le théorème de
Mihailescu~\cite{Mihailescu2004}, l'unique solution en entiers strictement
positifs à
\[
a^p - b^q = 1, \qquad a, b, p, q \geq 2,
\]
est $(a, b, p, q) = (3, 2, 2, 3)$. En spécialisant à notre cas
($p = 2$, $q \geq 3$), il vient $a = n+1 = 3$, $b = 2$, $q = 3$, et donc
$n = 2$ avec $n(n+2) = 2 \cdot 4 = 8 = 2^3$.

\medskip
En combinant les deux cas, $n(n+2) = b^q$ avec $b, q \geq 2$ impose
$n = 2$, $b = 2$, $q = 3$, et (\ref{eq:first-level}) devient l'identité
$3^2 - 2^3 = 1$. Réciproquement, $n = 2$ fournit évidemment une telle
représentation.
\end{proof}

\section*{3. Remarques}

\subsection*{3.1.\ Profondeur du théorème versus simplicité de l'application.}

Le théorème~\ref{thm:main} se réduit à une application en une ligne du
théorème de Mihailescu~\cite{Mihailescu2004}, mais le théorème sous-jacent
est profond~: sa preuve repose sur la théorie des unités cyclotomiques et
des racines primaires dans les corps cyclotomiques, en s'appuyant sur des
réductions antérieures dues à Cassels, Inkeri et d'autres. Ce qui est
nouveau ici n'est pas la profondeur, mais l'observation que cette
profondeur se manifeste dans un contexte entièrement élémentaire, celui
des radicaux emboîtés de Ramanujan. Le théorème est, en un sens, une
traduction de Mihailescu dans la famille de Ramanujan.

\subsection*{3.2.\ Niveaux plus profonds du radical.}

Le développement au premier niveau (\ref{eq:first-level}) est particulier.
À la profondeur $k$, des déploiements répétés donnent
\[
(n+1)^2 \;=\; 1 + n \, f(n+1)
\;=\; 1 + n\sqrt{1 + (n+1)\sqrt{1 + (n+2)\sqrt{\cdots}}},
\]
et les valeurs apparaissant à l'intérieur n'ont plus la structure de
$(n+1)^2 - 1$. Savoir si un niveau \emph{plus profond} du radical pour un
autre $n$ produit une identité de type Catalan est une question distincte~;
nous ne l'abordons pas ici, mais notons qu'elle exigerait une identité de
puissance parfaite non élémentaire dans
$\sqrt{1 + (n+1)\sqrt{\cdots}}$, dont aucune occurrence n'est connue.

\subsection*{3.3.\ Lien avec la conjecture de Pillai.}

La conjecture de Pillai affirme que pour tout $k \geq 1$, l'équation
$a^p - b^q = k$ admet un nombre fini de solutions $(a, b, p, q)$ avec
$a, b, p, q \geq 2$. Le cas $k = 1$ de Catalan est désormais établi
(Mihailescu)~; les $k$ supérieurs restent ouverts. Notre observation se
reformule ainsi~: \emph{au sein de la famille de Ramanujan}~(\ref{eq:Rama-general}),
la différence de puissances parfaites $|a^p - b^q|$ au premier niveau
atteint la valeur $1$ exactement une fois, en $n = 2$. Savoir si cela
reste vrai pour des écarts plus grands $|a^p - b^q| = k$, $k \geq 2$,
dépend des solutions paramétriques de $n(n+2) - b^q = \pm k$, où des
résultats partiels (par exemple Bilu, Bugeaud) s'appliquent, mais aucun
analogue général du théorème de Mihailescu n'est disponible.

\subsection*{3.4.\ Interprétation visuelle.}

Le théorème~\ref{thm:main} peut se paraphraser ainsi. La famille
(\ref{eq:Rama-general}) réalise tout entier strictement positif $n+1$ comme
radical emboîté~; parmi ceux-ci, l'entier $3$ est le seul dont le
\emph{produit définissant le plus interne} ($2 \cdot 4$) s'effondre
multiplicativement en une unique puissance parfaite ($2^3$). C'est ce qui
fait du cas $n = 2$, parmi tous les cas visuellement semblables de la
famille, un cas arithmétiquement singulier.

\begin{thebibliography}{9}

\bibitem{Berndt1985}
B.\ C.\ Berndt,
\emph{Ramanujan's Notebooks, Part~II},
Springer-Verlag, New York, 1989.

\bibitem{Herschfeld1935}
A.\ Herschfeld,
\emph{On infinite radicals},
Amer.\ Math.\ Monthly \textbf{42} (1935), no.\ 7, 419--429.

\bibitem{Mihailescu2004}
P.\ Mihailescu,
\emph{Primary cyclotomic units and a proof of Catalan's conjecture},
J.\ reine angew.\ Math.\ \textbf{572} (2004), 167--195.

\bibitem{Ramanujan1911}
S.\ Ramanujan,
\emph{Question 289},
J.\ Indian Math.\ Soc.\ \textbf{3} (1911), 90~;
\emph{Solution to Question 289},
J.\ Indian Math.\ Soc.\ \textbf{4} (1912), 226.

\bibitem{Vijayaraghavan1929}
T.\ Vijayaraghavan,
\emph{On certain infinite series and convergent integrals},
Proc.\ London Math.\ Soc.\ (2) \textbf{29} (1929), 281--294.

\end{thebibliography}

\end{document}
